\begin{figure}[ht!]

\centering
\small

\begin{tcolorbox}[
    boxsep=1pt,
    left=2pt,right=2pt,top=2pt,bottom=2pt,
    colback=blue!5!white,
    colframe=blue!40!black,
    before skip=2pt,
    after skip=2pt,
    sharp corners
]
\textbf{Question (Turn One):}
What is the IFRS 16 treatment of short term or low value leased assets? 
What are the restrictions?
\end{tcolorbox}

\begin{tcolorbox}[
    boxsep=1pt,
    left=2pt,right=2pt,top=2pt,bottom=2pt,
    colback=green!2!white,
    colframe=green!25!black,
    before skip=1pt,
    after skip=3pt,
    sharp corners
]
\small
\textbf{Gold Response:}
IFRS 16 Leases permit a simplified treatment for assets with a lease period of 12 months or less, or of low value. Although the standard does not give a numerical definition of 'low value' it does give examples of the types of assets that may be included, and this includes telephones. The simplified treatment allows the lease payments to be charged as an expense over the lease period, applying the accruals concept.

IFRS 16 Leases does not specify a particular monetary amount below which an asset would be considered 'low value', although the basis for conclusion indicates a value of \$5,000 as a guide. The standard also gives the following examples of low-value assets:

- tablets

- small personal computers

- telephones

- small items of furniture.

The assessment of whether an asset qualifies as having a 'low value' must be made based on its value when new. Therefore, a car would not qualify as a low-value asset, even if it was very old at the commencement of the lease.
\end{tcolorbox}

\begin{tcolorbox}[
    boxsep=1pt,
    left=2pt,right=2pt,top=2pt,bottom=2pt,
    colback=blue!5!white,
    colframe=blue!40!black,
    before skip=3pt,
    after skip=1pt,
    sharp corners
]
\small
\textbf{Question (Turn Two):}
On April 1st 2023, Abigail acquired telephones for her salesforce under a 
two-year lease agreement. The terms of the lease require an initial payment 
of \$3,000, followed by a payment of \$9,000 on March 31st 2024 and an 
additional \$9,000 on March 31st 2025. Show the impact of this lease 
arrangement on Abigail's financial statements for the year ended 
31 December 2023 under IFRS 16?
\end{tcolorbox}

\begin{tcolorbox}[
    boxsep=1pt,
    left=2pt,right=2pt,top=2pt,bottom=2pt,
    colback=green!2!white,
    colframe=green!25!black,
    before skip=1pt,
    after skip=2pt,
    sharp corners=downhill,
]
\small
\textbf{Gold Response:}
\[
\begin{aligned}
\text{Annual lease rental expense} &= \frac{\text{Total rentals payable}}{\text{Total lease period}} = \frac{\$3,000 + \$9,000 + \$9,000}{2\text{ years}} \\[6pt]
&= \$10,500\text{ per annum} \\[8pt]
\text{Expense to 31 December 2023} &= \$10,500 \times \frac{9}{12} = \$7,875 \\[8pt]
\text{Accrued expense} &= \$7,875 - \$3,000 = \$4,875
\end{aligned}
\]

The expense in this period of \$7,875 is not the same as the payment of \$3,000, so we need to accrue an additional expense of \$4,875.
\end{tcolorbox}


\caption{Untruncated example from BFF-Bench.}
\label{fig:bffbench-sample-full}
\end{figure}
